% I2MTC 2027 paper: Marvelmind vs OptiTrack on the KELO.
% Limit: 6 pages including references. Plan and numbers: ../idee.md.
% Build: needs IEEEtran.cls (in TeX Live / MiKTeX, or upload this folder to Overleaf).
%   pdflatex main && bibtex main && pdflatex main && pdflatex main
% Check the I2MTC 2027 author instructions (template, PDF eXpress) before submission.
\documentclass[conference]{IEEEtran}
\IEEEoverridecommandlockouts
\usepackage{cite}
\usepackage{amsmath,amssymb}
\usepackage{graphicx}
\usepackage{booktabs}
\usepackage{siunitx}
\usepackage{xcolor}
\graphicspath{{figures/}}

% \todo{...}: open points, remove all of them before submission.
\newcommand{\todo}[1]{\textcolor{red}{[TODO: #1]}}

\begin{document}

\title{Metrological Characterisation and Motion-Capture-Based Calibration of a Low-Cost Ultrasonic
Indoor Positioning System on a Mobile Robot\\
\todo{working title}}

\author{
\IEEEauthorblockN{Berat Baran Balci, \todo{co-authors}}
\IEEEauthorblockA{\textit{Department of Engineering, University of Ferrara}, Ferrara, Italy\\
\todo{e-mail addresses}}
}

\maketitle

% ---------------------------------------------------------------------------------------------------
\begin{abstract}
\todo{150--200 words, write last. One sentence each: problem (low-cost indoor positioning, reference
available only in part of the area); what was done (Marvelmind on the KELO vs OptiTrack, N runs);
main numbers (error at rest, while moving, delay); effect of the processing; result outside the
OptiTrack area.}
\end{abstract}

\begin{IEEEkeywords}
indoor positioning, ultrasonic positioning, motion capture, measurement uncertainty, mobile robots
\end{IEEEkeywords}

% ---------------------------------------------------------------------------------------------------
\section{Introduction}
% ~0.75 page
\todo{Low-cost indoor positioning systems (IPS) for mobile robots; why their accuracy must be
measured on the robot and not on a static stand.}

\todo{Motion capture gives sub-millimetre reference positions, but only in a small volume. Question:
how good is the IPS where the reference is available, and how well does a calibration made there
hold where it is not.}

Contributions:
\begin{itemize}
  \item \todo{characterisation of Marvelmind on the KELO against OptiTrack: at rest, while moving,
        vs speed, repeatability over repeated visits of the same points;}
  \item \todo{processing chain tuned with OptiTrack and validated with leave-one-run-out;}
  \item \todo{error outside the OptiTrack area (floor checkpoints, re-entry).}
\end{itemize}

% ---------------------------------------------------------------------------------------------------
\section{Related Work}
% ~0.5 page
\todo{Indoor positioning technologies and surveys~\cite{mautz2012indoor,zafari2019survey}.}
\todo{Published evaluations of Marvelmind and other ultrasonic IPS: search IEEE Xplore and add.}
\todo{Motion capture as reference and its own accuracy~\cite{merriaux2017vicon}.}
\todo{Trajectory error metrics (ATE)~\cite{sturm2012benchmark}; uncertainty
evaluation~\cite{gum2008}.}

% ---------------------------------------------------------------------------------------------------
\section{Experimental Setup}
% ~1 page
\subsection{Room and Marvelmind installation}
\todo{Room, 4 stationary beacons (heights, positions, quadrilateral area \SI{39.7}{m^2}), hedgehog
\#7, firmware 8.42, update rate \SI{5.76}{Hz}, filters off, temperature setting vs measured
temperature. The OptiTrack area is in a corner of the beacon quadrilateral: say it.}

\subsection{Reference: OptiTrack}
\todo{V120:Trio, \SI{120}{Hz}, Motive 3.5, covered area about \SI{2.7 x 2.7}{m}, calibration
residual, rigid body (markers, pivot), static noise at the stops.}

\subsection{Robot and mounting}
\todo{KELO driven with a joystick; hedgehog and marker positions; lever arm (measured and estimated).
Fig.~\ref{fig:setup}.}

\begin{figure}[t]
  \centering
  \fbox{\parbox[c][4cm][c]{0.9\columnwidth}{\centering \todo{photo of the robot + floor plan with
  beacons, OptiTrack area, triangle and floor checkpoints}}}
  % \includegraphics[width=\columnwidth]{setup.pdf}
  \caption{Experimental setup. \todo{caption}}
  \label{fig:setup}
\end{figure}

\subsection{Data logging and synchronisation}
\todo{Both systems logged on the same PC (\texttt{record.py}); Marvelmind device clock drift
(\SI{0.92}{\percent}) corrected; residual.}

% ---------------------------------------------------------------------------------------------------
\section{Method}
% ~1 page
\subsection{Frames and registration}
\todo{2D rigid transform (x, y, yaw) + z offset between OptiTrack and Marvelmind frames, lever arm
applied with the OptiTrack heading; estimated on calibration runs only~\cite{umeyama1991least}.
Note the mirrored Marvelmind map if not fixed.}

\subsection{Data selection}
\todo{Stop detection (OptiTrack speed below \SI{3}{cm/s} for at least \SI{10}{s}); exclusion of
frames with fallen or occluded markers, based on the Motive data.}

\subsection{Metrics}
\todo{At rest: per-stop offset and jitter. While moving: RMSE, 95th percentile, CDF, error vs speed,
effective delay. Outlier rate.}

\subsection{Processing chain}
\todo{(1) registration; (2) jump rejection; (3) temporal filter / delay compensation (Kalman,
constant velocity~\cite{kalman1960new}; offline smoother~\cite{rauch1965maximum}); (4) optional
beacon refinement. Leave-one-run-out.}

\subsection{Uncertainty budget}
\todo{OptiTrack, lever arm, synchronisation residual $\times$ speed, temperature
(speed of sound, about \SI{0.17}{\percent/\celsius})~\cite{gum2008}.}

% ---------------------------------------------------------------------------------------------------
\section{Results}
% ~2 pages
\subsection{At rest}
\todo{Offset and jitter at the stops; repeated visits of the triangle vertices: systematic vs
per-visit part.}

\subsection{While moving}
\todo{RMSE, p95, error vs speed, delay. Fig.~\ref{fig:speed}.}

\begin{figure}[t]
  \centering
  \fbox{\parbox[c][3.5cm][c]{0.9\columnwidth}{\centering \todo{error vs speed}}}
  \caption{Marvelmind error vs speed. \todo{caption}}
  \label{fig:speed}
\end{figure}

\subsection{Processing chain}
\todo{CDF for each step on the test runs; table.}

\begin{table}[t]
  \centering
  \caption{Error on the test runs (leave-one-run-out). \todo{fill}}
  \label{tab:summary}
  \begin{tabular}{lcccc}
    \toprule
    Step & Rest RMS & Moving RMS & Moving p95 & Delay \\
         & (cm) & (cm) & (cm) & (s) \\
    \midrule
    Raw (registration only) & -- & -- & -- & -- \\
    + jump rejection        & -- & -- & -- & -- \\
    + filter / delay comp.  & -- & -- & -- & -- \\
    \bottomrule
  \end{tabular}
\end{table}

\subsection{Outside the OptiTrack area}
\todo{Error on the floor checkpoints and at re-entry.}

% ---------------------------------------------------------------------------------------------------
\section{Discussion and Conclusion}
% ~0.5 page
\todo{Limits: one room, 2D only, joystick paths, fixed beacon geometry, OptiTrack area in a corner
of the beacon area. Main numbers in one paragraph.}

% ---------------------------------------------------------------------------------------------------
\bibliographystyle{IEEEtran}
\bibliography{refs}

\end{document}
